\section{Construction of the geodesic current}
\label{sec:construct}



\subsection{Countable additivity}
\label{sec:countable-additivity}

We prove countable sub-additivity of the function $\mu_{f}$ by
proving approximation results. We first find controlled sequences of
decreasing nested RH boxes that we can show have measure going to zero.

\begin{definition}
Let the RH lift system $(p_1,q_1,p_2,q_2)$ determine a right-handed
box
$B = [s_1,s_2) \times [t_2,t_1)$, where $s_i = \ell_i^-$ and $t_i =
r_i^+$. Further assume that $\ora{r_2}$ is nested inside $\ora{r_1}$,
so the box is of type~2 or~3.

Define the \emph{vertical accordion}, $F(p_1, q_1, p_2,\vec{q_2})$ of
boxes based at $B$ to be
the sequence of boxes, indexed by $k$, with $p_1, p_2, q_1$ constant
and $q_2$ translated by powers of $r_1 = q_1 x Q_1$:
\[
  p_1^{(k)}\coloneqq p_1,\quad
  p_2^{(k)}\coloneqq p_2,\quad
  q_1^{(k)}\coloneqq q_1,\quad
  q_2^{(k)} \coloneqq q_1 x^k q_1^{-1}q_2.
\]
This defines boxes whose  height becomes smaller as $k$
grows, with the top endpoint fixed. See
Figure~\ref{fig:vertical-accordion}.

\begin{figure}
\centering{
\resizebox{120mm}{!}{\svgfignormal{v-accordion}}
\caption{A vertical accordion $F(p_1,q_1, p_2,\vec{q_2})$, shown in
    the disk (left) and in the whiskers picture (right)}
  \label{fig:vertical-accordion}
}
\end{figure}

Similarly, if $\ora{\ell_1}$ is nested in $\ora{\ell_2}$ (so we again have a box of type 2 or 3), define the
\emph{horizontal accordion} $F(\vec{p_1},q_1, p_2,q_2)$ by
\[
  p_1^{(k)}\coloneqq p_1 P_2x^k p_2,\quad
  p_2^{(k)}\coloneqq p_2,\quad
  q_1^{(k)}\coloneqq q_1,\quad
  q_2^{(k)} \coloneqq q_1.
\]
\label{def:accordion}
\end{definition}
See Figure~\ref{fig:horizontal-accordion}.

\begin{proposition}
  Let $F = F(p_1,q_1,p_2,\vec{q_2})$ or $F(\vec{p_1},q_1, p_2,q_2)$ be a
  vertical or horizontal accordion as above. Then
  \begin{itemize}
  \item each $F(k)$ is a RH box;
  \item $F(k+1) \subset F(k)$; and
  \item $\displaystyle \lim_{k \to \infty} \wh{\mu_f}(F(k)) = 0$.
  \end{itemize}
\label{prop:finalfamily}
\end{proposition}

\begin{proof}
  We first prove the results in the case of vertical accordions,
  starting with checking that $F(k)$ is a box. Since $\ora{r_2}$ is
  nested inside $\ora{r_1}$, the NS dynamics of the group element $r_1
  = q_1 x q_1^{-1}$ translate $\ora{r_2}$ inside $(r_1^-,r_1^+)$, as depicted
  in Figure~\ref{fig:vertical-accordion}. In particular the cyclic
  order of the endpoints doesn't change, and we also automatically
  have $F(k+1) \subset F(k)$.

  To show the measures go to zero, we first reformulate the nesting
  condition algebraically. Note that
  $q_2Q_1 \cdot \ora{r_1} = \ora{r_2}$. Thus  $\ora{r_2}$ and $\ora{r_1}$ being nested means $(\ora{r_1},\ora{r_2})$ are R-parallel which is equivalent, by Lemma~\ref{lem:rhs}, to $(\ora{q_2Q_1}, \ora{r_1})$ R-crossing. This is
  in turn equivalent to $Q_2 \cdot \ora{q_2Q_1} = \ora{Q_1q_2}$
  R-crossing $\ora{x}$; write $e = Q_1q_2$.

  In the RH orbit notation, we then have
  \begin{align*}
    a^{(k)} &= p_1 q_1 = a &
    b^{(k)} &= p_1 q_2^{(k)} = p_1 q_1 x^k q_1^{-1}q_2 = a x^k e\\
    c^{(k)} &= p_2 q_1 = c &
    d^{(k)} &= p_2 q_2^{(k)} = p_2 q_1 x^k q_1^{-1} q_2 = c x^k e.
  \end{align*}

We thus get the family of RH orbits given by
\[
\bigl([B]_{a,a x^k e,c,c x^k e}\bigr)_{k=0}^\infty.
\]

To show this family of RH orbits have measure going to zero,
by definition of $\wh{\mu_f}$, we need to evaluate the nested limit:
\begin{align*}
  \lim_{k \to \infty} \bigl(\wh{\mu_f}([B]_{a,ax^ke,c,cx^ke})\bigr)
   &= \frac{1}{2}\lim_{k \to \infty} \lim_{n\to \infty} f([ax^ke x^n]) +
     f([cx^n])-f([ax^n])-f([cx^kex^n]) \\
  &= \frac{1}{2}\lim_{k,n \to \infty} \bigl(f([x^n a]) + f([x^ke])\bigr) + f([cx^n])\\
  &\qquad\qquad-f([ax^n])-\bigl(f([x^nc]) + f([x^ke])\bigr)\\
  &= 0.
\end{align*}
Here, we apply Lemma~\ref{lem:additivity-limit-2}
(adding and subtracting suitable copies of $f([x^n])$ and $f([x^k])$)
to see the double limit
exists, and then by applying Proposition~\ref{prop:join-split-alg} twice we
can break up the limit into canceling terms as desired.

The case of the horizontal accordion is similar, but there are some
confusing inverses so we spell it out briefly.
Since $P_2 p_1 \cdot \ora{\ell_1} = \ora{\ell_2}$,
the assumption that $\ora{\ell_1}$
is nested inside $\ora{\ell_2}$ (i.e., $(\ora{\ell_1},\ora{\ell_2})$ are R-parallel) is equivalent to $\ora{P_2 p_1}$ 
crossing $\ora{\ell_2}$ to the right, which in turn is equivalent to $\ora{p_1 P_2}$
crossing $\ora{x}$ to the right. In particular we have
\[
  \ora{\ell_1^{(k)}} = \bigl(p_1^{(k)}\bigr)^{-1} \cdot \ora{x}
  = \bigl(p_2^{-1} x p_2)^{-k} \cdot \ora{\ell_1}
  = \ell_2^{-k} \cdot \ora{\ell_1},
\]
which by NS dynamics of $\ell_2$ gives the picture on the left of
Figure~\ref{fig:horizontal-accordion}. In particular all of the $F(k)$
are boxes with $F(k+1) \subset F(k)$.

For the limiting computation, we set $h = p_1 P_2$, so that
\begin{align*}
  a^{(k)} &= p_1^{(k)} q_1 = h x^k c &
  b^{(k)} &= p_1^{(k)} q_2 = h x^k d\\
  c^{(k)} &= p_2 q_1 = c &
  d^{(k)} &= p_2 q_2 = d.
\end{align*}
Lemma~\ref{lem:additivity-limit-2} and
Proposition~\ref{prop:join-split-alg} apply as before.
\end{proof}

  \begin{figure}
\centering{
\resizebox{120mm}{!}{\svgfignormal{h-accordion}}
\caption{Horizontal accordion $F(\vec{p_1},q_1,p_2,q_2)$}
  \label{fig:horizontal-accordion}
}
\end{figure}

We now use Proposition~\ref{prop:finalfamily} to prove that $\mu_f$
is countably additive.
We will need to consider more types of boxes
than the half-open boxes we have been working with so far. For $B$ a
right-handed box $[s_1,s_2) \times [t_2,t_1)$, let
$\overline{B} \coloneqq [s_1,s_2] \times [t_2,t_1]$ and
$B^\circ \coloneqq (s_1,s_2) \times (t_2,t_1)$ be the corresponding
closed and open boxes.

\begin{figure}
     \centering
     \begin{subfigure}[b]{0.47\textwidth}
        \svgfignormal{ca1}
        \caption{$V(B,\epsilon)$ is a down and left extension of an RH
          box $B$ with mass less than $\epsilon$, as in Lemma~\ref{lem:extend1}}\label{fig:ca1}
     \end{subfigure}
     \hfill
     \begin{subfigure}[b]{0.47\textwidth}
        \svgfignormal{ca2}
        \caption{$W(B,\epsilon)$ is an $\epsilon$ neighborhood of the upper and right edges of an RH box $B$, as in Lemma~\ref{lem:extend2}}\label{fig:ca2}
     \end{subfigure}
     \caption{Extensions used in the proof of countable sub-additivity}
     \label{fig:cases}
\end{figure}

\begin{lemma}\label{lem:extend1}
  For $B = [s_1,s_2)\times [t_2,t_1) \in \cB$ a right-handed box and
  any $\epsilon > 0$, there is a right-handed box $V(B,\epsilon) =
  [s_1',s_2) \times [t_2',t_1) \in \cB$ so that
  $s_1' < s_1$ and $t_2' < t_2$ and
  $\wh{\mu_f}(V(B,\epsilon) \setminus B) < \epsilon$.
\end{lemma}
See Figure~\ref{fig:ca1}. That is, $V(B,\epsilon)$ extends $B$ slightly down and to the left,
without increasing the mass much.

\begin{proof}
  Let $B = B(p_1,q_1,p_2,q_2)$.
  By Proposition~\ref{prop:density}, there is some $q_3$ so that $q_3
  \cdot \ora{x}$ is nested inside $q_2 \cdot \ora{x}$; in particular
  $F_1 = F(p_1,q_2,p_2,\vec{q_3})$ is a vertical accordion, and for
  each $k$, $B \cup F_1(k)$ is a RH box. 
  By Proposition~\ref{prop:finalfamily}, there is $k_1$ large enough
  so that
  $\wh{\mu_f}(F_1(k_1))<\epsilon/2$.
  Repeat the procedure to find a horizontal
  accordion $F_2$ abutting on $B \cup F_1(k_1)$. Again there is $k_2$
  large enough so that $\wh{\mu_f}(F_2(k_2)) < \epsilon/2$, and set $V(B,\epsilon) \coloneqq B \cup F_1(k_1) \cup F_2(k_2)$.
\end{proof}

\begin{lemma}\label{lem:extend2}
  For $B = [s_1,s_2) \times [t_2,t_1) \in \cB$ a right-handed box and
  any $\epsilon > 0$, there is some $s_2',t_1'$ with
  $s_1 < s_2' < s_2$ and $t_2 < t_1' < t_1$ so that the region
  $W(B,\epsilon) = \bigl([s_1,s_2) \times [t_1',t_1)\bigr) \cup
  \bigl([s_2',s_2) \times [t_2,t_1)\bigr)$ is in $\mathcal{R}$ and satisfies
  $\wh{\mu_f}(W(B,\epsilon)) < \epsilon$.
\end{lemma}
See Figure~\ref{fig:ca2}. That is, $W(B,\epsilon)$ is a neighborhood of the upper and right
edges of~$B$ with small mass.

\begin{proof}
  Find $q_3$ so that $q_3 \cdot \ora{x}$ nested inside
  $q_1 \cdot \ora{x}$ and so that $F_1 = F(p_1,q_1,p_2,\vec{q_3})$ is a
  vertical accordion. Similarly find $p_3$ so that $P_3 \cdot \ora{x}$
  is nested inside $P_2 \cdot \ora{x}$ and
  $F_2 = F(\vec{p_3},q_1,p_2,q_2)$ is a horizontal accordion. Then
  for sufficiently large $k$, we have $\wh{\mu_f}(F_1(k)) < \epsilon/2$ and
  $\wh{\mu_f}(F_2(k)) < \epsilon/2$. Take
  $W(B,\epsilon) = F_1(k) \cup F_2(k)$.
\end{proof}
  
In the setting of Lemma~\ref{lem:extend2}, also set
\[
  W^\circ(B,\epsilon) \coloneqq \bigl([s_1,s_2] \times (t_1',t_1]\bigr) \cup
    \bigl((s_2',s_2] \times [t_1,t_2]\bigr);
\]
this is an open neighborhood of the upper and right edges of
$\overline{B}$ in the subspace topology of $\overline{B}$. Note that
neither $W(B,\epsilon)$ nor $W^\circ(B,\epsilon)$ is a subset of the
other.

\begin{proposition}\label{prop:countablyadditive}
  The finitely additive measure $\wh{\mu_f}$ is countably sub-additive on $\cB$.
\end{proposition}

\begin{proof}
  Let $B_0$ be a right-handed box, and suppose we are given a
  countable covering of it: $B_0 \subset \bigcup_{i=1}^\infty B_i$. We
  must show that
  \[
    \wh{\mu_f}(B_0) \le \sum_{i=1}^\infty \wh{\mu_f}(B_i).
  \]
  Fix $\epsilon > 0$. For $i \ge 1$, set
  $V_i \coloneqq V(B_i, \epsilon/2^i)$. Also add one more set by setting
  $V_0 \coloneqq W(B_0,\epsilon)$ and $V_0^\circ \coloneqq W^\circ(B_0,\epsilon)$.

  Then $\{V_i^\circ\}_{i=0}^\infty$ cover $\overline{B_0}$, since
  the $B_i$, for $i \geq 1$, cover $B_0$, $B_i \subset V_i^\circ$, and
  $\overline{B_0}\setminus B_0 \subset V_0^\circ$. Since
  $\overline{B_0}$ is compact and the $V_i^\circ \cap \overline{B_0}$ are open in the
  subspace topology, we have a finite subcover; let $I \subset
  \mathbb{N}$ be the set of
  indices included in the subcover, which necessarily contains $V_0^\circ$.

  By finite additivity and since $V_0^\circ \cap B_0 \subset V_0 \cap B_0$, we then have
  \begin{align*}
    \wh{\mu_f}(B_0) &\le \sum_{i \in I} \wh{\mu_f}(V_i)\\
      &\le \sum_{i \in I, i \ne 0} \wh{\mu_f}(B_i) + \sum_{i \in I, i \ne 0}
        \wh{\mu_f}(V_i \setminus B_i) + \wh{\mu_f}(V_0)\\
      &< \sum_{i=1}^\infty \wh{\mu_f}(B_i) + 2\epsilon.
  \end{align*}
  Since $\epsilon$ was arbitrary, we are done.
\end{proof}

We now put everything together to obtain an oriented geodesic current.

\begin{proposition}
Let $f \colon \Curves^+(S) \to \mathbb{R}$ be a functional on curves satisfying symmetry, additivity, smoothing and stability.
Then $\wh{\mu_f}$ extends to an oriented geodesic current $\mu_f$.
\end{proposition}
\begin{proof}
By Propositions \ref{prop:limitexists}, \ref{prop:RHsemiring},
\ref{prop:finitely-additive} and~\ref{prop:countablyadditive}, $\wh{\mu_f}$
is a finitely
additive and countably sub-additive pre-measure on $\mathcal{B}$, and by Proposition~\ref{prop:addsemitoring} has a unique finitely additive and countably sub-aditive extension to $\mathcal{R}=\mathcal{R}_{\mathcal{B}}$, which we also denote by $\wh{\mu_f}$. By covering $G(S)$ by a countable exhaustion consisting of a union of RH boxes, we see $\mu_f$ is $\sigma$-finite.
Therefore, by Carathéodory extension Theorem~\ref{thm:caratheodory},
$\wh{\mu_f}$ can be extend to a unique measure on
$\mathcal{A}_{\mathcal{R}}$, the $\sigma$-algebra generated
by~$\mathcal{R}$; we denote this extension by $\overline{\mu_f}$.
Let $\cA$ be the Borel $\sigma$-algebra of $G(S)$.
Since $\cR \subset \cA$, we have $\cA_{\cR} \subset \cA$.
On the other hand, by Proposition~\ref{prop:density}, the set of elements in
$\mathcal{R}$ is dense in the subspace of half-open rectangles of
$G(S)$. By the standard argument that half-rectangles with rational endpoints generate the Borel $\sigma$-algebra (see, e.g.,~\cite[Lemma~1.2.11]{Bo07:MeasureTheory}), it follows that $\cA \subset \cA_{\cR}$, hence $\overline{\mu_f}$ is a Borel measure.

Note that $\overline{\mu_f}$ is also locally finite, since any compact subset
of $G(S)$ can be covered by a finite union of the boxes in
$\mathcal{R}$, where the measure takes a finite value.
By definition of $\wh{\mu_f}$, the extension $\overline{\mu_f}$ is
$\pi_1(S)$-invariant. Finally, $\overline{\mu_f}$ is non-negative by
Proposition~\ref{prop:box-positive}.
\end{proof}

It is not \emph{a priori} clear from the definition of $\overline{\mu_f}$ that it is a flip-invariant geodesic current (recall
Definition~\ref{def:currents}).
In order to obtain a flip-invariant current, we can take
\[
\mu_f \coloneqq 1/2(\overline{\mu_f} + \sigma_*(\overline{\mu_f}))
\]
where $\sigma \colon G(S) \to G(S)$ is the flip map, and $\sigma_*$
the induced push-forward map on oriented geodesic currents.

\subsection{
Recovering \texorpdfstring{$f$}{f} on curves}
\label{sec:recovery}

It remains to show that $\mu_f$ recovers $f$ as an intersection number on curves.

If $y$ is a hyperbolic element representing a closed curve
$C$ in the surface $S$, choose
$z \in \partial_{\infty} S \setminus \{ y_-, y_+\}$ so that
$y^+,y^-,\, z,\, y \cdot z$ are in cyclic order.
Then
$(y^+,y^-) \times [z,y \cdot z)$ is a Borel fundamental domain for the
  action of $\langle y \rangle$ on the set of geodesics that intersect
  $\ora{y}$ transversally.

Before proceeding, we discuss an asymmetric version of intersection number between oriented geodesic currents.

\begin{definition} The \emph{asymmetric intersection} number between oriented geodesic currents is defined as follows
\begin{equation}
\vec{i}(\mu,\nu) \coloneqq \int_{\mathcal{F}} \mathbf{1}_R \, d(\mu \otimes \nu),
\label{eq:asymmetric}
\end{equation}
where $\mathcal{F}$ is a fundamental domain of the diagonal action of $\pi_1(S)$ on $\DG(S)$, given by the fact 
where the function $\mathbf{1}_R \colon \DG(S) \to \mathbb{R}$ is defined as
\[
\mathbf{1}_R(l,m) \coloneqq
\begin{cases}
1 & \text{if } (l,m)\text{ are R-crossing},\\
0 & \text{otherwise}.
\end{cases}
\]
\label{def:asymmetric}
\end{definition}
From the definition, we have the following relation between the standard intersection form between oriented geodesic currents and the asymmetric one.
\begin{equation}
i(\mu, \nu) = \vec{i}(\mu, \nu) + \vec{i}(\mu, \sigma(\nu)).
\label{eq:relation_intersections}
\end{equation}

For $p,q \in \wt{S}$, one can write  $\ora{G[p,q)}\subset G(S)$ for a \emph{right-oriented} transversal, i.e., the subset of
oriented geodesics in $G[p,q)$ which right cross the bi-infinite $\ora{l}$ oriented geodesic obtained by extending $[p,q) \subset \wt{S}$. Similarly, we will use $\ole{G[p,q)}\subset G(S)$ for the \emph{left-oriented} transversal, i.e., the subset of
oriented geodesics in $G[p,q)$ intersecting the inverted geodesic $\ora{L}$ to the right. Recall that given a box $B$, we denote by $\ora{B}$ its associated right-oriented box, and we can similarly define a left-oriented box $\ole{B}$.
We have the following partitions.
\begin{subequations}\label{eq:partition_oriented}
\begin{align}
G[p,q) &= \ora{G[p,q)} \cup \ole{G[p,q)} 
\label{eq:partition_transversal_oriented}\\
B &= \ora{B} \cup \ole{B}.
\label{eq:partition_box_oriented}
\end{align}
\end{subequations}

We can compute intersections with oriented curves as follows.

\begin{lemma}The intersection and asymmetric intersection
of a closed curve $C=[y]$ with an oriented geodesic current $\mu$ can be computed as 
\[
i(\mu, C) = \mu([z,y \cdot z) \times (y^+,y^-))
\]
and
\[
\vec{i}(\mu, C) = \mu(\ora{[z,y \cdot z) \times (y^+,y^-)}).
\]
\end{lemma}
\begin{proof}
The second claim follows from the observation that, for any $q \in \ora{y}$, $\ora{G[q, y\cdot q)}$ and $\ora{[z, y\cdot z) \times (y^+, y^-)}$ are two fundamental domains for the action of $\langle y \rangle$ on the set of geodesics that intersect $\ora{y}$ to the right.
Specifically, if $y \in \pi_1(S)$ denotes the element representing the oriented curve $C$, then 
\[
\mu(\ora{G[q, y \cdot q)})= \mu(\ora{[z,y \cdot z) \times (y^+,y^-)})
\]
for every $q \in \ora{y}$ and for every $z \notin \{ y^-, y^+\}$ so that $z, y\cdot z, y^+, y^-$ are in the cylic order. See~\cite[Lemma~4.8]{HP97:RigidityNegCurvedCone}): the proof only uses
the $\pi_1(S)$-invariance of $\mu$. 
The same claim follows for the intersection number using Equation~\ref{eq:relation_intersections}.
\end{proof}



We finish by proving Theorem~\ref{thm:intersect}, first in the following
form.

\begin{theorem}
Let $f \colon \Curves^+(S) \to \mathbb{R}$ be a curve functional satisfying symmetry, smoothing, additivity and stability, and $\mu_f, \overline{\mu_f}$ as above.
For $C$ a closed curve, we have
\[
i(\mu_f,C)=\vec{i}(\overline{\mu_f},C)=f(C).
\]
In particular, $\mu_f$ is unique among flip-invariant geodesic currents.
\label{prop:recoverintersection}
\end{theorem}
\begin{proof}
We first prove that $\vec{i}(\overline{\mu_f},C)=f(C)$. Given a representative $y \in
\pi_1(S)$ of the conjugacy class~$C$, we will find a sequence of
RH boxes to approximating the fundamental domain for $\langle
y\rangle$ above, and then show that computing the measure of that
fundamental domain yields $f([y])$.

Choose $p,q \in \pi_1(S)$ so that $(P\cdot \ora{x}, \ora{y},
q\cdot\ora{x})$ are R-parallel, as in Figure~\ref{fig:recover-f}.
The NS dynamics of the action of $y$ imply that, for $m_1
> m_2$ and $n_1 > n_2$, the box
$B(py^{-m_1},py^{-m_2},y^{n_1}q,y^{n_2}q)$ form a RH box. We will look
at the following particular family for $n \ge 0$:
\begin{align*}
  p_1^{(n)} &\coloneqq p & q_1^{(n)}&\coloneqq y^n q\\
  p_2^{(n)} &\coloneqq p y &q_2^{(n)} &\coloneqq y^{-(n+1)} q.
\end{align*}

\begin{figure}
\centering{
\svgfignormal{recover-f}}
\caption{The boxes used in proving that $\mu_f$ recovers $f$.}
  \label{fig:recover-f}
\end{figure}

\begin{comment}
%Hand drawn figure
\begin{figure}
\centering{
\includegraphics[height=88mm]{recover-f}
\caption{The boxes used in proving that $\mu_f$ recovers $f$.}
\label{fig:recover-f}
}
\end{figure}
\end{comment}
In the RH orbit notation, we have
\begin{align*}
  a^{(n)} &= p y^{n}   q &  b^{(n)} &= p y^{-(n+1)} q \\
  c^{(n)} &= p y^{n+1} q &  d^{(n)} &= p y^{-n} q.
\end{align*}

If we let $z=(P \cdot \ora{x})^-$, and $B_n \coloneqq B_{a^{(n)}, b^{(n)}, c^{(n)}, d^{(n)}}$ then we have
    \begin{multline*}
      \vec{i}(\mu, [y])=\overline{\mu_f}(\ora{[z,Y \cdot z))\times (y^+,y^-)}) \\
      \begin{aligned}
    &= \lim_{n\to\infty} \overline{\mu_f}(B_n)  \\
    &= \frac{1}{2}\lim_{n\to\infty}\lim_{m\to\infty}f([py^{n+1}qx^{m}]) + f([py^{-(n+1)}qx^{m}])-f([py^{n}qx^{m}]) - f([py^{-n}qx^{m}]) \\
    &= \frac{1}{2}\bigl(f([y]) + f([y^{-1}])\bigr) =f([y]).
      \end{aligned}
    \end{multline*}
    The first equality follows by the fact that the boxes defined
    above converge from below to the  $\ora{[z,Y
      \cdot z) \times (y^+,y^-)}$ (the fundamental domain for the
  action of $\langle Y \rangle$) as $n \to \infty$, since $B_{n-1} \subset B_n$.
    At the second equality, we use definition of $\wh{\mu_f}$ on RH boxes.
    For the third, Lemma~\ref{lem:additivity-limit-2} lets us pass
    from the iterated limit to a double limit, and gives a
    number of canceling terms of $mf([x])$ and $nf([y^{\pm 1}])$; only
    two terms survive the cancellation. Then symmetry of $f$ gives the result.

      The property $i(\mu_f, C)=f(C)$ for any closed curve $C$ persists by the definition of $\mu_f$ and the fact that $\overline{\mu_f}$ is finitely additive. Therefore, $\mu_f$ is a geodesic current with the same intersection numbers as the oriented geodesic current $\overline{\mu_f}$.

 Finally, the uniqueness of the geodesic current $\mu_f$ is
 a result of Otal, geodesic currents are determined by their marked
 length spectrum \cite[Th\'eor\`eme~2]{Otal90:SpectreMarqueNegative}.
\end{proof}

\begin{remark}
In \cite[Proposition~4.3]{BIPP24:PositiveCrossratios} a geodesic current is induced
from a cross-ratio.
Our construction is more general (in particular, we do not rely on a cross-ratio).
We discuss here why our construction does not follow from their arguments.
Given a dense subset $X \subset \partial_{\infty} S$, they consider the set of $\mathcal{R}(X)$ of all proper rectangles of geodesics with endpoints on $X$.
They start from a finitely additive function $r \colon \mathcal{R}(X) \to \mathbb{R}$ (defined suitably from a cross-ratio)
that is (1) side-by-side additive: if $R,R_1,R_2 \in \mathcal{R}(X)$
with $R=R_1 \cup R_2$, with $R_1,R_2$ disjoint, $r(R)=r(R_1)+r(R_2)$;
and (2) monotonic: if $R,R' \in \mathcal{R}(X)$ with $R \subset R'$, we have $r(R) \leq r(R')$.

Our function $\mu_f$ on RH boxes satisfies these properties by Proposition~\ref{prop:sidebyside}
and Proposition~\ref{prop:finitely-additive}, however, if we take $X$ to
be the set of
endpoints of translates of~$x$, then $\mu_f$ is only
defined in a proper subset of $\mathcal{R}(X)$
Therefore, we cannot directly apply their result.

Moreover, we need the extra control on the family of RH
boxes to recover the curve functional~$f$ from the geodesic current.
In Section~\ref{sec:crossratios}, we prove that every cross-ratio induces a curve functional (i.e., period) satisfying the hypotheses of our Theorem~\ref{thm:intersect}, and thus our result generalizes \cite[Proposition~4.3]{BIPP24:PositiveCrossratios}.
\end{remark}




\subsection{Proof of  Theorem~\texorpdfstring{\ref{thm:intersect}}{A}}

Now we prove the theorem as stated in the introduction: in terms of functionals on unoriented multi-curves.

\begin{proof}[Proof of Theorem~\ref{thm:intersect}]

By Definition~\ref{def:smoothing} and Proposition~\ref{prop:power_disconnected_imply_stable}, the induced additive and symmetric functional on oriented multi-curves satisfies stability. Hence, by Theorem~\ref{prop:recoverintersection}, the result follows.
\end{proof}

\subsection{Proof of Theorem~\texorpdfstring{\ref{thm:intersections_group}}{A'}}
\label{subsec:algebraic_vs_top}
We finish this section by proving the algebraic version of the
theorem, namely, Theorem~\ref{thm:intersections_group}.

\begin{proof}[Proof of Theorem~\ref{thm:intersections_group}]
As in the statement, let $\overline{f}(a) \coloneqq
\frac{1}{2}\left(f(a) + f(A) \right)$, and extend it linearly to be an
additive
function on multi-curves.

To prove the theorem, we will show that (this extension of) $\overline{f}$ satisfies the hypothesis of Theorem~\ref{thm:intersect}, and then apply that Theorem to conclude.


We will first verify smoothing of $\overline{f}$.
We start
with oriented disconnected smoothing, between crossing elements
$(\ora{a},\ora{b})$, which we can assume R-cross by switching $a,b$ if
necessary. We have
\begin{align*}
  \overline{f}(a) + \overline{f}(b)
    &= \frac{1}{2}\bigl(f(a) + f(A) + f(b) + f(B)\bigr)\\[2pt]
    &\ge \frac{1}{2}\bigl(f(ab) + f(AB)\bigr) = \overline{f}(ab).
\end{align*}
In the inequality, we used the assumption on $f$ and the fact that
$(\ora{A},\ora{B})$ also
R-cross, as well as conjugation invariance for $f(AB) = f((ba)^{-1})$.
By grouping the four terms differently, and using the fact that
$(\ora{b},\ora{A})$ R-cross, we also get $\overline{f}(a) +
\overline{f}(b) \ge \overline{f}(bA)$.

Next we turn to oriented connected smoothing of a curve at an
essential self-crossing.
Using Proposition~\ref{prop:crossings-triality}, assume the essential
self-crossing is realized by the self-intersection of a curve
$[\gamma]=[ab]$, where $(\ora{a},\ora{b})$ are parallel; by switching
them if necessary, we can assume they are R-parallel, and thus so are
$(\ora{B},\ora{A})$. Then
\begin{align*}
  \overline{f}(ab)
    &= \frac{1}{2}\bigl(f(ab) + f(BA)\bigr) \\[2pt]
    &\ge \frac12\bigl(f(a) + f(b) + f(B) + f(A)\bigr)
     = \overline{f}(a) + \overline{f}(b).
\end{align*}

We have thus deduced $\overline{f}$ satisfies oriented connected and disconnected smoothing. Therefore, by Corollary~\ref{cor:oriented_smoothing_unoriented} below,
$\overline{f}$~satisfies unoriented connected and disconnected smoothing as well. 
Finally, stability implies oriented power-smoothing, i.e., $f(a^{n+m}) \geq f(a^n) + f(b^n)$ for every $n, m \geq 1$. This establishes smoothing of $\overline{f}$. Hence, smoothing, stability and additivity are satisfied for the function on unoriented curves $\overline{f}$.
We are thus under the
hypotheses of Theorem~\ref{thm:intersect}.
\end{proof}

\begin{remark}
In fact,  we can replace the following hypotheses in~Theorem~\ref{thm:intersections_group} by the following slightly weaker ones:
\begin{itemize}
 \item $f$ can be replaced by power smoothing (by Proposition~\ref{prop:power_disconnected_imply_stable});
\item oriented smoothing can be replaced by R-oriented smoothing:
 \begin{align}
f(ab) &\ge f(a) + f(b)
&& \text{if } (\ora{a},\ora{b}) \text{ are R-parallel}, \\
f(a) + f(b) &\ge f(ab)
&& \text{if } (\ora{a},\ora{b}) \text{ R-cross}.
\end{align}
\end{itemize}
\end{remark}



%%% Local Variables:
%%% mode: latex
%%% TeX-master: "Intersections"
%%% End:
